\documentclass[11pt,a4paper]{article}

% ---------------------------------------------------------------
%  Publication List -- Stanisław Szufa
%  EPFL IC Faculty Search 2026
%  The entries come from publications.bib via make_publist.py.
% ---------------------------------------------------------------
\usepackage[T1]{fontenc}
\usepackage[utf8]{inputenc}
\usepackage{mathpazo}
\renewcommand{\ttdefault}{lmtt}
\usepackage[margin=2.5cm]{geometry}
\usepackage{microtype}
\usepackage{enumitem}
\usepackage{titlesec}
\usepackage{xcolor}
\usepackage[hidelinks]{hyperref}
\hypersetup{pdfauthor={Stanisław Szufa},pdftitle={Stanisław Szufa: Publication List}}

\definecolor{accent}{RGB}{160,20,30}   % same accent as the other statements
\titleformat{\section}{\large\bfseries\color{accent}}{}{0pt}{}
\titlespacing*{\section}{0pt}{10pt}{4pt}
\setlength{\parindent}{0pt}
\setlength{\parskip}{4pt}
\setlist{nosep,leftmargin=1.2em}
\interlinepenalty=10000  % never split a publication entry across pages

% ---------------------------------------------------------------
%  Shared document header -- Stanisław Szufa, EPFL application
%  Load after xcolor and hyperref, and after defining the accent colour:
%    % ---------------------------------------------------------------
%  Shared document header -- Stanisław Szufa, EPFL application
%  Load after xcolor and hyperref, and after defining the accent colour:
%    % ---------------------------------------------------------------
%  Shared document header -- Stanisław Szufa, EPFL application
%  Load after xcolor and hyperref, and after defining the accent colour:
%    \input{../shared/header}
%  Usage:
%    \docheader{Research Statement}        % name, title, contact line, rule
%    \docheader{}                          % no title (cover letter)
%    \docheader[extra contact line]{Title} % second contact line (CV)
%    \headercontactfalse before \docheader    % no contact lines
% ---------------------------------------------------------------
\providecolor{muted}{RGB}{90,90,90}
\newcommand{\headersep}{\enspace{\color{accent!60}\textbullet}\enspace}
% Contact lines are shown unless a document says \headercontactfalse.
\newif\ifheadercontact\headercontacttrue
\newcommand{\docheader}[2][]{%
  \begin{center}
    {\LARGE\bfseries\color{accent} Stanisław Szufa}\par
    \if\relax\detokenize{#2}\relax\else
      \vspace{3pt}{\large\color{muted} #2}\par
    \fi
    \ifheadercontact
    \vspace{5pt}
    {\small\color{muted}
    University of Geneva, Switzerland\par\vspace{1pt}
    \href{mailto:stanislaw.szufa@unige.ch}{stanislaw.szufa@unige.ch}\headersep
    +41 76 472 16 09%
    \if\relax\detokenize{#1}\relax\else\headersep#1\fi\par}
    \fi
  \end{center}
  \vspace{-6pt}
  {\color{accent!40}\rule{\linewidth}{0.5pt}}\par
  \vspace{2pt}}

%  Usage:
%    \docheader{Research Statement}        % name, title, contact line, rule
%    \docheader{}                          % no title (cover letter)
%    \docheader[extra contact line]{Title} % second contact line (CV)
%    \headercontactfalse before \docheader    % no contact lines
% ---------------------------------------------------------------
\providecolor{muted}{RGB}{90,90,90}
\newcommand{\headersep}{\enspace{\color{accent!60}\textbullet}\enspace}
% Contact lines are shown unless a document says \headercontactfalse.
\newif\ifheadercontact\headercontacttrue
\newcommand{\docheader}[2][]{%
  \begin{center}
    {\LARGE\bfseries\color{accent} Stanisław Szufa}\par
    \if\relax\detokenize{#2}\relax\else
      \vspace{3pt}{\large\color{muted} #2}\par
    \fi
    \ifheadercontact
    \vspace{5pt}
    {\small\color{muted}
    University of Geneva, Switzerland\par\vspace{1pt}
    \href{mailto:stanislaw.szufa@unige.ch}{stanislaw.szufa@unige.ch}\headersep
    +41 76 472 16 09%
    \if\relax\detokenize{#1}\relax\else\headersep#1\fi\par}
    \fi
  \end{center}
  \vspace{-6pt}
  {\color{accent!40}\rule{\linewidth}{0.5pt}}\par
  \vspace{2pt}}

%  Usage:
%    \docheader{Research Statement}        % name, title, contact line, rule
%    \docheader{}                          % no title (cover letter)
%    \docheader[extra contact line]{Title} % second contact line (CV)
%    \headercontactfalse before \docheader    % no contact lines
% ---------------------------------------------------------------
\providecolor{muted}{RGB}{90,90,90}
\newcommand{\headersep}{\enspace{\color{accent!60}\textbullet}\enspace}
% Contact lines are shown unless a document says \headercontactfalse.
\newif\ifheadercontact\headercontacttrue
\newcommand{\docheader}[2][]{%
  \begin{center}
    {\LARGE\bfseries\color{accent} Stanisław Szufa}\par
    \if\relax\detokenize{#2}\relax\else
      \vspace{3pt}{\large\color{muted} #2}\par
    \fi
    \ifheadercontact
    \vspace{5pt}
    {\small\color{muted}
    University of Geneva, Switzerland\par\vspace{1pt}
    \href{mailto:stanislaw.szufa@unige.ch}{stanislaw.szufa@unige.ch}\headersep
    +41 76 472 16 09%
    \if\relax\detokenize{#1}\relax\else\headersep#1\fi\par}
    \fi
  \end{center}
  \vspace{-6pt}
  {\color{accent!40}\rule{\linewidth}{0.5pt}}\par
  \vspace{2pt}}


\begin{document}

\headercontactfalse
\docheader{Publication List}

\textbf{Summary.} 5 journal articles and 28 conference papers,
published mainly at the major AI conferences 
(AAAI, IJCAI, NeurIPS, AAMAS) and in
\emph{Artificial Intelligence}, \emph{JAIR}, and \emph{JCSS}.

\vspace{6pt}

\textbf{Author ordering.} As is customary in computational social choice,
authors are usually listed alphabetically, so position in the author list
does not reflect contribution. Papers that deviate from alphabetical order
list the lead author first.

\vspace{4pt}

% Generated by make_publist.py from publications.bib -- do not edit by hand.

\section{Journal Articles}
\begin{itemize}[label={},leftmargin=3.2em,itemsep=2.5pt]
\item[\textbf{[J5]}] P. Skowron, A. Slinko, \textbf{S. Szufa}, and N. Talmon. Participatory Budgeting with Cumulative Votes. \emph{Theory and Decision}, 2025.
\item[\textbf{[J4]}] \textbf{S. Szufa}, N. Boehmer, R. Bredereck, P. Faliszewski, R. Niedermeier, P. Skowron, A. Slinko, and N. Talmon. Drawing a Map of Elections. \emph{Artificial Intelligence} 343, 2025. {\small\color{accent}Journal version of AAMAS-2020 and IJCAI-2021 papers.}
\item[\textbf{[J3]}] P. Faliszewski, P. Skowron, A. Slinko, K. Sornat, \textbf{S. Szufa}, and N. Talmon. How Similar Are Two Elections? \emph{Journal of Computer and System Sciences} 150, 2025. {\small\color{accent}Journal version of the AAAI-2019 paper.}
\item[\textbf{[J2]}] P. Faliszewski, K. Sornat, and \textbf{S. Szufa}. The Complexity of Subelection Isomorphism Problems. \emph{Journal of Artificial Intelligence Research} 80, pp.~1343--1371, 2024. {\small\color{accent}Journal version of the AAAI-2022 paper.}
\item[\textbf{[J1]}] N. Boehmer, K. Heeger, and \textbf{S. Szufa}. A Map of Diverse Synthetic Stable Matching Instances. \emph{Journal of Artificial Intelligence Research} 79, pp.~1113--1166, 2024. {\small\color{accent}Journal version of the AAMAS-2023 paper.}
\end{itemize}

\section{Conference Papers}
\begin{itemize}[label={},leftmargin=3.2em,itemsep=2.5pt]
\item[\textbf{[C28]}] R. Faulkner, \textbf{S. Szufa}, D. Hoyer, T. J. Zhang, R. Bouffanais, J. Z. Leibo, and Z. Jin. Simulating Democratic Deliberation: Representing Pluralistic Preferences through Electoral Systems in Multi-Agent {LLMs}. \emph{Findings of EMNLP-2026}, pp.~11673--11705.
\item[\textbf{[C27]}] P. Faliszewski, J. Mertlov{\'a}, K. Sornat, \textbf{S. Szufa}, and T. Wąs. Agreement, Diversity, and Polarization Indices for Approval Elections. \emph{Proceedings of IJCAI-2026}, pp.~3393--3401.
\item[\textbf{[C26]}] P. Faliszewski, K. Sornat, \textbf{S. Szufa}, and T. Wąs. Outer Diversity of Structured Domains. \emph{Proceedings of AAMAS-2026}, pp.~2106--2114.
\item[\textbf{[C25]}] P. Böhm, R. Bredereck, P. Gölz, A. Kaczmarczyk, and \textbf{S. Szufa}. Putting Fair Division on the Map. \emph{Proceedings of AAAI-2026}, pp.~16726--16734.
\item[\textbf{[C24]}] P. Faliszewski, K. Sornat, \textbf{S. Szufa}, and T. Wąs. Diversity of Structured Domains via k-Kemeny Scores. \emph{Proceedings of AAAI-2026}, pp.~16880--16888.
\item[\textbf{[C23]}] P. Faliszewski, Ł. Janeczko, A. Kaczmarczyk, G. Lisowski, P. Skowron, \textbf{S. Szufa}, and M. Szwagierczak. Strategic Cost Selection in Participatory Budgeting. \emph{Proceedings of NeurIPS-2025}.
\item[\textbf{[C22]}] F. Nikolow, P. Faliszewski, and \textbf{S. Szufa}. Maps of Tournaments: Distances, Experiments, and Data. \emph{Proceedings of ECAI-2025}, pp.~675--682.
\item[\textbf{[C21]}] D. Boratyn, W. S{\l}omczy{\'n}ski, D. Stolicki, and \textbf{S. Szufa}. Spoiler Susceptibility in Party Elections. \emph{Proceedings of ECAI-2025}, pp.~3775--3782.
\item[\textbf{[C20]}] P. Faliszewski, Ł. Janeczko, A. Kaczmarczyk, M. Kurdziel, G. Pierczyński, and \textbf{S. Szufa}. Learning Real-Life Approval Elections. \emph{Proceedings of AAMAS-2025}, pp.~704--712.
\item[\textbf{[C19]}] P. Faliszewski, J. Mertlov{\'{a}}, P. Nunn, \textbf{S. Szufa}, and T. Wąs. Distances Between Top-Truncated Elections of Different Sizes. \emph{Proceedings of AAAI-2025}, pp.~13823--13830.
\item[\textbf{[C18]}] T. Delemazure, Ł. Janeczko, A. Kaczmarczyk, and \textbf{S. Szufa}. Selecting the Most Conflicting Pair of Candidates. \emph{Proceedings of IJCAI-2024}, pp.~2766--2773.
\item[\textbf{[C17]}] N. Boehmer, P. Faliszewski, Ł. Janeczko, D. Peters, G. Pierczyński, Š. Schierreich, P. Skowron, and \textbf{S. Szufa}. Evaluation of Project Performance in Participatory Budgeting. \emph{Proceedings of IJCAI-2024}, pp.~2678--2686.
\item[\textbf{[C16]}] W. Słomczyński, D. Stolicki, and \textbf{S. Szufa}. Nonparametric Detection of Gerrymandering in Multiparty Plurality Elections. \emph{Proceedings of IJCAI-2024}, pp.~2967--2975.
\item[\textbf{[C15]}] N. Boehmer, P. Faliszewski, Ł. Janeczko, A. Kaczmarczyk, G. Lisowski, G. Pierczyński, S. Rey, D. Stolicki, \textbf{S. Szufa}, and T. Wąs. Guide to Numerical Experiments on Elections in Computational Social Choice. \emph{Proceedings of IJCAI-2024}, pp.~7962--7970. {\small\color{accent}Survey track.}
\item[\textbf{[C14]}] G. Pierczyński and \textbf{S. Szufa}. Single-Winner Voting with Alliances: Avoiding the Spoiler Effect. \emph{Proceedings of AAMAS-2024}, pp.~1567--1575.
\item[\textbf{[C13]}] Ł. Janeczko, J. Lang, G. Lisowski, and \textbf{S. Szufa}. Discovering Consistent Subelections. \emph{Proceedings of AAMAS-2024}, pp.~935--943.
\item[\textbf{[C12]}] P. Faliszewski, M. Lackner, K. Sornat, and \textbf{S. Szufa}. An Experimental Comparison of Multiwinner Voting Rules on Approval Elections. \emph{Proceedings of IJCAI-2023}, pp.~2675--2683.
\item[\textbf{[C11]}] P. Faliszewski, J. Flis, D. Peters, G. Pierczyński, P. Skowron, D. Stolicki, \textbf{S. Szufa}, and N. Talmon. Participatory Budgeting: Data, Tools and Analysis. \emph{Proceedings of IJCAI-2023}, pp.~2667--2674.
\item[\textbf{[C10]}] P. Faliszewski, A. Kaczmarczyk, K. Sornat, \textbf{S. Szufa}, and T. Wąs. Diversity, Agreement, and Polarization in Elections. \emph{Proceedings of IJCAI-2023}, pp.~2684--2692.
\item[\textbf{[C9]}] N. Boehmer, K. Heeger, and \textbf{S. Szufa}. A Map of Diverse Synthetic Stable Roommates Instances. \emph{Proceedings of AAMAS-2023}, pp.~1003--1011. {\small\color{accent}Pragnesh Jay Modi Best Student Paper Award.}
\item[\textbf{[C8]}] N. Boehmer, R. Bredereck, E. Elkind, P. Faliszewski, and \textbf{S. Szufa}. Expected Frequency Matrices of Elections: Computation, Geometry, and Preference Learning. \emph{Proceedings of NeurIPS-2022}.
\item[\textbf{[C7]}] \textbf{S. Szufa}, P. Faliszewski, Ł. Janeczko, M. Lackner, A. Slinko, K. Sornat, and N. Talmon. How to Sample Approval Elections? \emph{Proceedings of IJCAI-2022}, pp.~496--502.
\item[\textbf{[C6]}] N. Boehmer, P. Faliszewski, R. Niedermeier, \textbf{S. Szufa}, and T. Wąs. Understanding Distance Measures Among Elections. \emph{Proceedings of IJCAI-2022}, pp.~102--108.
\item[\textbf{[C5]}] P. Faliszewski, K. Sornat, and \textbf{S. Szufa}. The Complexity of Subelection Isomorphism Problems. \emph{Proceedings of AAAI-2022}, pp.~4991--4998.
\item[\textbf{[C4]}] N. Boehmer, R. Bredereck, P. Faliszewski, R. Niedermeier, and \textbf{S. Szufa}. Putting a Compass on the Map of Elections. \emph{Proceedings of IJCAI-2021}, pp.~59--65.
\item[\textbf{[C3]}] \textbf{S. Szufa}, P. Faliszewski, P. Skowron, A. Slinko, and N. Talmon. Drawing a Map of Elections in the Space of Statistical Cultures. \emph{Proceedings of AAMAS-2020}, pp.~1341--1349.
\item[\textbf{[C2]}] P. Faliszewski, P. Skowron, A. Slinko, \textbf{S. Szufa}, and N. Talmon. How Similar Are Two Elections? \emph{Proceedings of AAAI-2019}, pp.~1909--1916.
\item[\textbf{[C1]}] P. Faliszewski, \textbf{S. Szufa}, and N. Talmon. Optimization-Based Voting Rule Design: The Closer to Utopia the Better. \emph{Proceedings of AAMAS-2018}, pp.~32--40.
\end{itemize}

\section{Preprints and Manuscripts}
\begin{itemize}[label={},leftmargin=3.2em,itemsep=2.5pt]
\item[\textbf{[P2]}] \textbf{S. Szufa}, D. Scherer, D. Stolicki, R. Bouffanais, and D. Wernli. Abstention Patterns in the United Nations General Assembly. \emph{Manuscript}, 2026.
\item[\textbf{[P1]}] U. Grandi, R. Bredereck, T. Delemazure, U. Endriss, J. Maly, N. Mattei, N. Maudet, O. Nardi, and \textbf{S. Szufa}. Social Choice with Text: Collective Decision Making in the {LLM} Era. \emph{Position paper, preprint, HAL hal-05548117}, 2026. {\small\color{accent}Arising from a working group at Dagstuhl Seminar 25401.}
\end{itemize}


% ===============================================================
\section{Thesis}
% ===============================================================
\begin{itemize}[label={},leftmargin=3.2em,itemsep=2.5pt]
  \item[\textbf{[T1]}] \textbf{S. Szufa}. \emph{Map of Elections}. PhD thesis,
  AGH University of Kraków, 2024. Defended with distinction; Polish Prime
  Minister's Award for an outstanding doctoral dissertation.
\end{itemize}

% ===============================================================
\section{Software and Datasets}
% ===============================================================
\begin{itemize}[label={},leftmargin=3.2em,itemsep=2.5pt]
  \item[\textbf{[S1]}] \textbf{Pabulib} (\href{https://pabulib.org}{pabulib.org}).
  Open library of participatory budgeting elections; 2{,}053 instances from 30
  cities in 8 countries, cited in more than 100 publications. Founded in
  2020.
  \item[\textbf{[S2]}] \textbf{prefsampling}. Python library for
  sampling synthetic preference data under a wide range of statistical
  cultures. Co-developed with Simon Rey.
  \item[\textbf{[S3]}] \textbf{mapof}. Python library for
  computing distances between elections and drawing maps of elections.
\end{itemize}

\end{document}
